\chapter{RE-014 ---  Design Doctrine}\label{record-400}

\noindent{\footnotesize\color{BookMuted}\path{docs/method/retro/RE_014-buffer-holds-facts-doctrine.md}\par}\smallskip

\begin{RecordMeta}
title: RE-014 — "The Buffer Holds Facts" Design Doctrine
lane: retro
legend: RE

\end{RecordMeta}

\section{Disposition}\label{record-400-disposition}

Implemented on \texttt{r\allowbreak{}e\allowbreak{}l\allowbreak{}e\allowbreak{}a\allowbreak{}s\allowbreak{}e\allowbreak{}/\allowbreak{}v\allowbreak{}5\allowbreak{}.\allowbreak{}0\allowbreak{}.\allowbreak{}0} as an explicit invariant. The new \href{https://github.com/flyingrobots/bijou/blob/cbd70612140e7a450d764b3c7be1c6af878b99bc/docs/method/invariants/buffer-holds-facts.md}{\texttt{d\allowbreak{}o\allowbreak{}c\allowbreak{}s\allowbreak{}/\allowbreak{}i\allowbreak{}n\allowbreak{}v\allowbreak{}a\allowbreak{}r\allowbreak{}i\allowbreak{}a\allowbreak{}n\allowbreak{}t\allowbreak{}s\allowbreak{}/\allowbreak{}b\allowbreak{}u\allowbreak{}f\allowbreak{}f\allowbreak{}e\allowbreak{}r\allowbreak{}-\allowbreak{}h\allowbreak{}o\allowbreak{}l\allowbreak{}d\allowbreak{}s\allowbreak{}-\allowbreak{}f\allowbreak{}a\allowbreak{}c\allowbreak{}t\allowbreak{}s\allowbreak{}.\allowbreak{}m\allowbreak{}d}} captures the rule that queues, buffers, and message channels should carry facts rather than executable behavior, and the architecture signposts now point at that invariant from the root, core, and TUI runtime docs.

\section{Original Proposal}\label{record-400-original-proposal}

Legend: \href{https://github.com/flyingrobots/bijou/blob/1853ad902f6bf219326409e7f743a5a1d6517dbc/docs/method/legends/RE-runtime-engine.md}{RE --- Runtime Engine}

\section{Idea}\label{record-400-idea}

During the RE-007 buffer migration, James articulated a principle:

\begin{quote}
``Keep buffered shell commands as plain discriminated data and interpret them explicitly. Use a handler table if you want to avoid a giant switch. The buffer should hold facts, not code.''
\end{quote}

This is a design doctrine that should be captured in the runtime engine's architectural documentation. It applies beyond shell commands --- any buffer, queue, or event channel in the system should carry data (what happened) not behavior (how to respond).

